\documentclass[runningheads]{llncs}

\usepackage[T1]{fontenc}
\usepackage{graphicx}

\begin{document}

\title{ArtBench-10 Generative Modeling: VAE, DCGAN, and Diffusion}
\author{Nuno Batista \and Miguel Cabral Pinto}
\institute{University of Coimbra}
\maketitle

\begin{abstract}
This report presents the design, training, and evaluation of three generative
model families on ArtBench-10: a VAE, a DCGAN, and diffusion-based models. As well as some variants of the diffusion model and the GAN, including a pretrained Google DDPM fine-tuning approach. We detail our data pipeline, model architectures, training protocols, and evaluation metrics (FID/KID).
We compare model behavior on the provided 20\% subset, select the strongest
approach, and discuss expected full-dataset retraining.

\end{abstract}

\section{Introduction}
% Motivation for generative modeling in artistic domains.
% Project scope and expected contributions.


\section{Problem Description}
% ArtBench-10 characteristics (32x32 RGB, 10 styles, 50k train).
% Why this is a generation task and not classification.
% Main challenges: fidelity, diversity, style consistency, stability.

\section{Methodology and Approach}

\subsection{Data Pipeline}
% Data source and folder structure used in the repository.
% Normalization, batching, subset protocol (20% CSV subset).
% Full-data retraining plan.

\subsection{Model Families}

\subsubsection{VAE}
% Architecture summary, latent dimension, beta, reconstruction objective.

\subsubsection{DCGAN}
% Generator/discriminator design, spectral norm usage, adversarial objective.

\subsubsection{Diffusion (From Scratch)}
% UNet design, noise schedule, training objective, sampling strategy.

\subsubsection{Diffusion (Pretrained Google DDPM Fine-Tuning)}
% Pretrained checkpoint and rationale.
% Freezing strategy (down/mid frozen, up trainable).
% Low learning rate and scheduler locking.
% Unconditional setup and observed metric/perception trade-offs.

\subsection{Implementation and Reproducibility}
% Config-driven experiments (YAML), random seeds, checkpointing.
% W\&B logging policy and artifact tracking.
% Hardware/software setup.

\section{Experimental Setup}

\subsection{Training Protocol}
% Hyperparameter tables per model.
% Number of epochs, optimizer settings, batch size.

\subsection{Evaluation Protocol}
% Mandatory protocol from statement:
% - 5000 generated samples
% - 5000 real samples
% - FID on full sets
% - KID with 50 subsets of size 100
% - 10 repetitions with different seeds
% Report mean +- std for FID/KID.

\section{Results and Analysis}

\subsection{Qualitative Results}
% Include sample grids for each model family.
% Comment on fidelity/diversity/style coherence.

\subsection{Quantitative Results}
% Table: FID mean +- std, KID mean +- std (all required runs).
% Keep evaluation code path fixed across models.

\subsection{Ablations and Discussion}
% Explain major design choices and observed effects.
% Example: pretrained diffusion visual quality vs FID/KID behavior.
% Example: flat diffusion noise-MSE during fine-tuning.

\section{Conclusion}
% Final comparison summary.
% Best model on subset and rationale.
% Full-data retraining status/final outcome.
% Limitations and short future work notes.

\section*{Use of External Resources and AI Tools}
% List external repositories, pretrained components, and AI assistants used.
% Describe how each was used and how outputs were validated.

\begin{thebibliography}{10}
% Add project references here (ArtBench, VAE, DCGAN, DDPM, evaluation metrics, etc.).
\end{thebibliography}

\end{document}
